\scp{Target audience for this study}{01-02}
We are certainly hoping that experienced Austronesian scholars will benefit from this study,
but they are not our primary target audience for the organization
and discussion in this book.
Some have been reading
and thinking for decades about Austronesian languages,
about subgrouping arguments,
about the methodology,
the debates in the literature,
challenges to working with old sources,
and even about related disciplines of pre-history
and archaeology, anthropology, and DNA studies.
They may already be familiar to varying degrees with some of the data
and studies and arguments relating to our target region.

There are issues in our target region
which have taken us years to grow into an understanding of them.
And there are things that each of us has shared with the other that have opened up new insights.
Having regularly worked
and corresponded with other developing
and established scholars,
both of us are aware of things that are often assumed,
things that are not known,
and things that are helpful to know about the Austronesian languages in our target region.
Hence, in some places we elaborate on topics
that more senior Austronesian linguists might feel “it goes without saying”.
Yet our interactions with established
and emerging scholars indicate that some things need such “unnecessary” elaboration.
The target audience for this study includes:

\begin{itemize}
	\item Austronesian linguists who may not be familiar with the target region,
				its literature, data, and complexities.
	\item Investigators new to the region,
				be they comparative linguists,
				descriptive linguists, anthropologists, archaeologists, or other.
				Consequently, we try to lay out the data
				and explain the arguments for people who are not comparative linguists to also be able to see the patterns we have found.
	\item Indonesian and East Timorese scholars who are interested in the prehistory
				and classification of the languages within our target region,
				one or more of which they might speak,
				but perhaps have had limited access to the scattered modern resources that try to deal with the broader topic,
				or lay out the data in ways that they can also see
				and understand the patterns
				and the issues.
				This volume brings many scattered issues together in one place.
\end{itemize}

This target audience is why we use the word “unravelling” in the book title.
In addition to the primary linguistic data,
there are issues in related disciplines that provide helpful background
and orientation for the region as the context in which the languages are spoken.
There are various methodological traps in the transcription
and glosses of old sources that,
when alerted to them,
can save much time
and confusion and error in one’s own conclusions.
A number of these issues have been misunderstood in the comparative literature,
and consequently have contributed to confusion regarding the languages under investigation in this book.

We want to share what we have learned the hard way.
There is no need for everyone else to have to rediscover (or fail to discover) some of the complexities
and pitfalls that we have encountered along the way.